\chapter{GUI, Distributed Execution, and Running the Flow}
\label{ch:gui}

\glance{Drive OpenROAD interactively with the Qt GUI (\pkg{src/gui}) or the
browser viewer (\pkg{src/web}); scale heavy stages with the distributed
framework (\pkg{src/dst}). The autonomous RTL-to-GDSII wrapper is
\textbf{OpenROAD-flow-scripts} (ORFS). This chapter ties the earlier stages
into a runnable flow and a command cheat-sheet.}

\section{GUI: \pkg{src/gui}}

\id{gui::Gui} is the Qt layout viewer. Launch with \cmd{openroad -gui} or
\cmd{gui::show} from a Tcl session. It renders the odb design (instances,
wires, special nets), exposes layer visibility, heatmaps (density, power, IR
drop, congestion), selection/inspection, and a Tcl console --- every flow
command works from inside the GUI. A \id{Renderer} plugin API lets modules
draw overlays (grt/drt/est register this way).

\section{Web viewer: \pkg{src/web}}

\pkg{src/web} serves the layout as PNG tiles over WebSocket for browser-based
zoom/pan without a native display:

\begin{lstlisting}[style=ortcl]
web_server -port 8080 -dir /path/to/web/assets
save_image -web out.png          ;# headless-friendly
\end{lstlisting}

\section{Distributed compute: \pkg{src/dst}}

\id{dst::Distributed} provides a load-balancer / worker TCP framework used to
parallelize compute-heavy stages (notably detailed routing's
\cmd{-distributed} mode). Workers and a balancer are started as separate
OpenROAD processes and share a volume for intermediate data.

\section{RAM generator: \pkg{src/ram}}

\id{ram::RamGen} (experimental, DFFRAM-inspired) builds dense stdcell-based
memory arrays when a foundry compiler is unavailable:

\begin{lstlisting}[style=ortcl]
generate_ram -width 32 -depth 128 -ports 1rw
\end{lstlisting}

Supports 1rw / 1rw1r / 1r1w / 2r1w; tested on sky130hd and Nangate45.

\section{Running the OpenROAD binary}

\begin{lstlisting}[style=ortcl]
openroad                     # interactive Tcl
openroad -gui                # with Qt GUI
openroad -python script.py   # Python entry
openroad -exit flow.tcl      # batch, exit when done
openroad -metrics out.json -threads 8 flow.tcl
\end{lstlisting}

Example scripts live under \file{test/} (e.g.\ \file{gcd\_nangate45.tcl},
\file{aes\_sky130hd.tcl}) and share a common \file{flow.tcl}.

\section{OpenROAD-flow-scripts (ORFS)}

ORFS (sibling repo \file{openroad\_repos/OpenROAD-flow-scripts}) is the
Makefile-driven autonomous RTL-to-GDSII wrapper:

\begin{enumerate}
  \item \textbf{Yosys} synthesizes RTL $\rightarrow$ gate netlist (often
        calling ABC).
  \item \textbf{OpenROAD} runs the P\&R stages from earlier chapters.
  \item Per-platform / per-design \file{config.mk} files under
        \file{flow/platforms} and \file{flow/designs} set PDK, SDC, utilization,
        etc.
\end{enumerate}

ORFS is where ``one command builds a chip''; lives; OpenROAD is the engine it
calls.

\section{Ecosystem (siblings in \file{openroad\_repos/})}

\begin{center}
\small
\begin{tabularx}{\linewidth}{@{}lX@{}}
\toprule
\textbf{Repo} & \textbf{Role} \\
\midrule
OpenSTA & standalone timer (also vendored as \pkg{src/sta}) \\
OpenROAD-MCP & Model Context Protocol server for AI tooling \\
ORAssistant & RAG assistant over OpenROAD docs \\
abc & Berkeley ABC (also embedded at \file{src/abc}; see ABC manual) \\
\bottomrule
\end{tabularx}
\end{center}

\section{Command cheat-sheet}

\begin{center}
\small
\begin{tabularx}{\linewidth}{@{}lXX@{}}
\toprule
\textbf{Stage} & \textbf{Primary command(s)} & \textbf{Ch.} \\
\midrule
Read design     & \cmd{read\_lef}, \cmd{read\_verilog}, \cmd{link\_design}, \cmd{read\_liberty} & \ref{ch:odb}, \ref{ch:timing} \\
Floorplan       & \cmd{initialize\_floorplan}, \cmd{make\_tracks}, \cmd{tapcell} & \ref{ch:floorplan} \\
I/O pins        & \cmd{place\_pins} & \ref{ch:place} \\
Macros          & \cmd{rtl\_macro\_placer} & \ref{ch:place} \\
Global place    & \cmd{global\_placement} & \ref{ch:place} \\
Legalize        & \cmd{detailed\_placement}, \cmd{filler\_placement} & \ref{ch:place} \\
PDN             & \cmd{define\_pdn\_grid} \ldots \cmd{pdngen} & \ref{ch:power} \\
Pre-CTS repair  & \cmd{estimate\_parasitics -placement}, \cmd{repair\_design}, \cmd{repair\_timing} & \ref{ch:rsz} \\
CTS             & \cmd{clock\_tree\_synthesis} & \ref{ch:cts} \\
Global route    & \cmd{global\_route} & \ref{ch:route} \\
Post-GR repair  & \cmd{estimate\_parasitics -global\_routing}, \cmd{repair\_timing} & \ref{ch:rsz} \\
Detailed route  & \cmd{detailed\_route} & \ref{ch:route} \\
Extract / fill  & \cmd{extract\_parasitics}, \cmd{density\_fill} & \ref{ch:timing}, \ref{ch:route} \\
IR drop         & \cmd{analyze\_power\_grid} & \ref{ch:power} \\
Write out       & \cmd{write\_def}, \cmd{write\_db} & \ref{ch:odb} \\
\bottomrule
\end{tabularx}
\end{center}

\glance{To understand OpenROAD: master OpenDB + STA (Ch.~\ref{ch:odb},
\ref{ch:timing}), then walk this cheat-sheet top to bottom. Every later
module is ``read the DB, compute, write the DB.''}
